% CECI EST UN FICHIER DE MODÈLE LATEX POUR LES ARTICLES INCLUS DANS
% *Anthology of Computers and the Humanities*. AJOUTEZ L'OPTION
% 'final' LORS DE LA CRÉATION DE LA VERSION FINALE DE L'ARTICLE. 
% NE PAS modifier la classe du document
%\documentclass[final,fra]{anthology-ch} % pour la version finale
\documentclass[final,fra]{anthology-ch}         % pour la soumission

% CHARGER LES PACKAGES LaTeX
\usepackage{booktabs}
\usepackage{graphicx}

% AJOUTEZ vos propres packages avec \usepackage{}

% TITRE DE LA SOUMISSION
% Remplacez ceci par le titre de votre soumission
\title{La puissance décolonisatrice des affects complexes}

% INFORMATIONS SUR LES AUTEURS ET LES AFFILIATIONS
% Pour chaque auteur, utilisez une nouvelle commande \author, avec
% les numéros entre crochets indiquant les affiliations correspondantes
% (section suivante) et l'ORCID-ID pour chaque auteur.  
\author[1]{Lia Beatriz Torraca}[
  orcid=0000-0002-3860-8500,
  email=metaffectus@gmail.coml
]

\affiliation{1}{Universidade do Estado do Rio de Janeiro, Brésil}

% MOTS-CLÉS
% Fournissez un ou plusieurs mots-clés séparés par des virgules
% en utilisant les commandes suivantes
%\keywords[fra]{informatique, sciences humaines}
\keywords[eng]{affect, algorithms, data colonialism, photography, Instagram}

\keywords[fra]{affect, algorithmes, colonialisme des données, photographie, Instagram}

% MÉTADONNÉES DE LA PUBLICATION
% Ces champs seront remplis lors de la publication; 
% ils peuvent être laissés par défaut lors de la soumission
\pubyear{2026}
\pubvolume{4}
\pagestart{1}
\pageend{7}
\paperorder{1}
\conferencename{Actes de la Conférence Humanistica}
\conferenceeditors{Serena Crespi, Simon Gabay, Martin Grandjean, Ariane Pinche, Marie Puren et Léa Saint-Raymond}
\doi{00000/00000}
\customhead{}

\addbibresource{bibliography.bib}

%%%%%%%%%%%%%%%%%%%%%%%%%%%%%%%%%%%%%%%%%%%%%%%%%%%%%%%%%%%%%%%%%%%%%%%%%%%
% DÉBUT DU TEXTE
\begin{document}

\maketitle

\begin{abstract}
    Although virtualized life opens up new perspectives for the human imagination, it is also responsible for creating new forms of vulnerability, especially among teenagers and young adults who are social media users. Faced with the challenges posed by the algorithmic governance imposed by platforms such as Instagram, the project ``Portray, Reset and Reterritorialize: decolonizing the algorithmically configured digital experience through affect'' aims to offer strategies that enable the decolonization of these users' experiences. Strategies developed from observing the role of complex affect during photographic experiences generated and mediated by algorithms. On this occasion, we will present the main theoretical and methodological aspects applied, which constitute a proposal for digital visual education through photography and new forms of image production. This research is anchored in Silvan Tomkins' Theory of Affects and Maurice Merleau-Ponty's Phenomenology of Perception.
\end{abstract}

\section{Introduction}\label{introduction-1}

Ce projet découle d'une réflexion sur la gouvernance esthétique à laquelle nous sommes soumis. Le projet aborde l'actualisation de l'esthétique colonialiste, imposée par les plateformes de réseaux sociaux selon le concept d'extractivisme numérique de Nick Couldry et Ulises Méjias \cite{couldry2019costs}. Les auteurs ont nommé ce phénomène \textit{data colonialism} \cite{couldry2019costs}, responsable de la configuration esthétique des utilisateurs de plateformes digitales comme \textit{Instagram}. Cette configuration esthétique signifie la modulation de la perception de l´utilisateur de ces plateformes et le calibrage de son organisation affective à des niveaux profonds. L'interférence s'étend à la compréhension du monde et de l'existence par le sujet dans une nouvelle dimension spatio-temporelle, correspondant à la déterritorialisation des corps et des objets, comme Lévy comprend le virtuel \cite{levy2007virtual}. Cette interférence a été exacerbée par la pandémie de COVID-19, lorsque nous avons commencé à ressentir intensément les processus de virtualisation.  

Dans ce contexte, le projet a été conçu pour construire des stratégies capables de faire face aux défis posés par ce nouveau régime esthétique ce qui représente le phénomène du colonialisme des données, notamment pour les adolescents (15-17 ans) et jeunes adultes (18-32 ans), qui représentent une part importante des utilisateurs des plateformes de réseaux sociaux, selon le rapport \textit{Data Not Found: Social Media Data Transparency for Information Integrity} \cite{santini_data_2026}. Il s'agit de les protéger des menaces qu'impliquent un environnement visuel contrôlé par les Big Techs, qui, dans ce projet de recherche, fait référence à Meta Platforms Inc., propriétaire de la plateforme Instagram.

À cet égard, le projet vise à comprendre l'expérience de ce public sur la plateforme Instagram. Cela implique d'observer des affects qui émergent dans ces relations communicatives et dans les expériences photographiques programmées prévues dans le cadre du projet. Cette observation peut contribuer au développement de stratégies pour une littératie numérique élargie, qui permettra la décolonisation de l´expérience digitale de ces utilisateurs, au-delà de la plateforme Instagram.

Le cadre théorique s'articule autour d'études sur l'affect, la perception et l'image, ancrées dans la théorie des affects de Silvan Tomkins \cite{tomkins1962affect,tomkins1963affect}, associée à la phénoménologie de la perception de Maurice Merleau-Ponty \cite{merleauponty2018fenomenologia}, comme référence méthodologique. Une recherche qualitative, descriptive et prépositive, étayée par une discussion théorique et des pratiques psychoéducatives avec activités photographiques structurées en \textit{Focus Groups} \cite{abreu2009grupos} avec des jeunes. Un projet qui pourrait permettre une réflexion sur ce nouveau modèle de pouvoir imposé aux utilisateurs d'Instagram.

Notre recherche se fonde sur le concept d'affect complexe développé par Torraca \cite{torraca2019espetaculo}. Un concept inspiré par la pensée d'Edgar Morin \cite{morin2015introducao} et conçu par l'association de deux autres concepts: l'affect neutre de Silvan Tomkins \cite{tomkins1962affect,tomkins1963affect} et le virtuel de Pierre Lévy \cite{levy2007virtual}. L'affect complexe est un affect neutre virtualisé.

\section{La puissance décolonisatrice des affects complexes}

%Le projet a été structuré sur le concept d´affect complexe \cite{torraca2024antropofagia}, développé par des études sur la théorie des affects de Silvan Tomkins \cite{tomkins1962affect,tomkins1963affect} et le concept de virtuel de Pierre Lévy \cite{levy2007virtual}.
Le couplage d'affects complexes permet au sujet de réinitier son expérience digitale à partir des expériences photographiques programmées. Cette programmation photographique fait partie d'une proposition pour l'éducation visuelle numérique.

La théorie des affects de Tomkins \cite{tomkins1962affect,tomkins1963affect} distingue les affects primaires, considérés comme innés, en affects positifs, affects négatifs et un seul affect neutre: l'affect de surprise. Ces affects primaires constituent le système biologique sous-jacent à l'émotion et participent à l'acte de percevoir, contribuant à moduler les émotions et à produire un répertoire imagé \cite{tomkins1962affect}. Les affects primaires peuvent influencer les relations sociales et, par conséquent, l'environnement et les affects sociaux dominants, qui se reflètent dans le système communicationnel. Pour Tomkins, l'affect neutre a une capacité particulière: réinitialiser les expériences du sujet, c'est pourquoi Tomkins l'a appelé \emph{resetting affect} \cite{tomkins1962affect}. Selon Torraca \cite{torraca2024antropofagia}, cet affect de réinitialisation permettrait de briser les stéréotypes et de formuler d'autres images lorsqu'il est associé à l'expérience esthétique programmée comme l'expérience photographique. L'affect de surprise est l'élément introduit dans l'expérience du sujet déterritorialisé, agissant comme un agent de décolonisation, capable d'intervenir dans sa configuration affective et perceptive, lui permettant ainsi de réinitialiser son expérience esthétique.

Selon Lévy \cite{levy2007virtual}, les émotions complexes sont liées à la problématisation dans le processus de déterritorialisation, ainsi qu'à la génération de nouvelles modalités affectives provoquées par le transit entre le réel et le virtuel \cite{levy2007virtual}. Pour l'auteur, le sujet situé dans le monde virtuel est un sujet affectif, et sa perception peut être altérée, ce qui se refléterait dans sa configuration esthétique, plus précisément dans son organisation perceptive et affective. Les affects seraient responsables de l'actualisation du monde virtuel, provoquant l'émergence de nouveaux types d'affects, que l'on pourrait considérer comme une inventivité affective. L'affect complexe est ce nouveau type d'affect: un affect neutre déterritorialisé.

Sous l'effet d'algorithmes, soit médiateurs soit génératifs, l'affect complexe peut assumer deux modalités: l'affect virtuel et l'affect artificiel \cite{torraca2024antropofagia}. Le premier émerge de la médiation algorithmique des plateformes de réseaux sociaux, agissant comme un affect médiateur des expériences de l'utilisateur dans ces espaces. Le second est généré par la \og communication artificielle\fg{} \cite{esposito2022artificial} établie entre les utilisateurs et les algorithmes génératifs, comme dans le cas des expériences du \emph{photo-prompt} (modalité d´expérience photographique), agissant comme un affect \emph{réinitialisateur}. Autrement dit, l'affect artificiel permet à l'utilisateur de réinitialiser son expérience esthétique. Pour Torraca \cite{torraca2024antropofagia}, le partage d'images générées par des algorithmes génératifs sur les plateformes de réseaux sociaux permet le couplage d'affects complexes \cite{torraca2024antropofagia}. Ce couplage affectif permet la réécriture du \textit{script} algorithmique configuré à partir de la modification des images et la recréation du monde perçu, ce qui traduit les processus de décolonisation. Il s'agit de la possibilité d'une émancipation du sujet, lui permettant de promouvoir sa propre (re)configuration esthétique, telle que la reterritorialisation des relations sociales, qui se reflète au-delà des plateformes de réseaux sociaux.

\section{Le Projet}

Le projet vise à élaborer des stratégies à travers de nouvelles formes de production d'images, en essayant des expériences photographiques virtuelles dans des espaces contrôlés tels qu'Instagram. Ce sont des stratégies fondées sur la connaissance et la compréhension des émotions qui émergent de ces expériences. Parmi ces stratégies, nous utilisons des pratiques psychoéducatives avec des activités photographiques structurées en groupes de discussion \cite{abreu2009grupos} avec des jeunes. Selon Ponciano \cite{ponciano2024praticas}, ces pratiques contribuent à 
\begin{itemize}
    \item réduire les inégalités d'accès à l'information;
    \item mettre en œuvre des actions de prévention;
    \item renforcer la capacité à gérer le stress et les problèmes de santé mentale;
    \item à favoriser l'autonomie, l'entraide et des environnements propices au bien-être.
\end{itemize}
 \noindent Les pratiques psychoéducatives \cite{ponciano2024praticas} prennent en compte les dimensions émotionnelles et relationnelles. Sur la base de ces expériences programmées et de l'observation des affects qui émergent dans ces groupes de discussion \cite{abreu2009grupos}, il est possible de déclencher des processus de décolonisation et de construire des stratégies d'émancipation capables de promouvoir l'éducation au visuel numérique.

Face aux enjeux posés par la gouvernance algorithmique, le projet a été structuré autour de l'étude de l'expérience des jeunes utilisateurs d'Instagram, et plus particulièrement des affects qui en émergent. La compréhension de cette expérience peut étayer le développement de stratégies qui permettent une littératie élargie, réunissant les dimensions numériques, médiatiques, algorithmiques, affectives et juridiques. Cette littératie élargie est essentielle pour 
%les utilisateurs 
comprendre les processus affectifs et perceptifs dans un monde radicalement axé sur l'image. Ces stratégies permettent de stimuler la créativité de l'utilisateur sans compromettre la fluidité de sa communication. Ainsi, nous pensons qu'il est possible pour l'utilisateur d'intervenir dans le scénario algorithmique défini par Instagram. Nous proposons donc des stratégies que les utilisateurs d'Instagram peuvent élaborer, fondées sur le couplage d'affects complexes \cite{torraca2024antropofagia}. Ces stratégies peuvent être développées grâce à l'expérience photographique et à de nouveaux modèles de production d'images.

La photographie, considérée dans toute l'étendue de ses possibilités, y compris dans le domaine de la post-photographie
\footnote{Selon Joanna Zylinska \cite{zylinska_perception_2023}, professeure de philosophie des médias et de pratiques numériques critiques au King's College de Londres, le terme \og post-photographie\fg{} s'est imposé au cours de la dernière décennie pour décrire la nouvelle réalité du médium photographique. Ce terme a émergé du débat sur les similitudes et les différences entre la photographie analogique et numérique, ainsi que sur le rôle des \og amateurs\fg{} dans la production photographique. Le terme \og post-photographie\fg{} a gagné en visibilité en 2011, suite à un manifeste de Joan Fontcuberta, publié dans le journal \textit{La Vanguardia}, qui annonçait une nouvelle ère pour la photographie, qui a laissé dans le passé certaines caractéristiques de ce médium dans la construction de la réalité et de la mémoire. Un terme qui acquiert une nouvelle pertinence avec les systèmes algorithmiques génératifs et les nouvelles formes de production d'images, comme ce qu'on appelle déjà la du \og Photo Prompt\fg{} ou \og Photographie du Prompt\fg{}. Il est important de noter que Zylinska utilise le terme \textit{After Photography} [14], proche de la dimension affective, liée à l'idée fondamentale de notre projet.}
, est le \emph{medium} qui permet de coupler des affects complexes \cite{torraca2024antropofagia}. Nous pensons que l'utilisation de l'intelligence artificielle générative dans les activités psychoéducatives \cite{ponciano2024praticas} peut contribuer à la pensée critique. Cette dynamique permettrait l'introduction de la littératie digitale, que nous considérons comme essentielle à l'émancipation de ce sujet en transformation.

Dans ce projet, nous avons choisi de combiner deux approches méthodologiques pour le développement de la recherche: la méthode phénoménologique merleau-pontyenne \cite{merleauponty2018fenomenologia} et la théorie des affects de Silvan Tomkins \cite{tomkins1963affect,tomkins1962affect}, liée aux manifestations des affects. La méthode de Merleau-Ponty permet d'élargir la perspective du sujet, en lui offrant des voies de reconfiguration esthétique, c'est-à-dire la possibilité de modifier sa perception et son organisation affective d'après ce que comprend Torraca \cite{torraca2019espetaculo}. L'une de ces voies est la synthèse perceptive, à partir de laquelle le sujet commence à sensibiliser son regard et à développer ses capacités esthétiques.

La théorie des affects de Silvan Tomkins \cite{tomkins1963affect,tomkins1962affect} s'approche de la phénoménologie de Merleau-Ponty \cite{merleauponty2018fenomenologia}, car tous deux considèrent l'expérience comme une source d'observation, embrassant toute sa complexité. La théorie des affects se présente comme un fondement théorique et méthodologique pour l'observation et l'étude des affects qui émergent lors d'une expérience esthétique déterritorialisée et artificiellement configurée. Cette théorie prend en compte la complexité des affects générés et médiatisés algorithmiquement, centrés sur le corps, même si celui-ci se situe dans un autre espace-temps. À partir de la méthode phénoménologique, il est possible d'explorer de nouvelles formes pour décoloniser les affects et les relations établies sur les plateformes avec un régime esthétique similaire à celui d'Instagram, considérant que ce réseau social constitue un espace radicalement imagé.

En complément de cette approche, nous utilisons la technique du \emph{Focus Group} pour la recherche qualitative \cite{abreu2009grupos}. Les \textit{Focus Groups} sont des groupes de discussion dont la programmation est basée sur des expériences photographiques. Cette technique s'avère être une stratégie méthodologique pertinente pour parvenir à une compréhension approfondie des manifestations affectives dans les expériences photographiques proposées \cite{abreu2009grupos}. Elle permet aussi d'explorer les aspects connexes de ce nouveau régime esthétique colonialiste imposé aux utilisateurs d'Instagram. C'est une technique qui favorise la spontanéité et le naturel des échanges, tout en élargissant le champ d'observation des manifestations affectives liées à l'image et générées par l'interaction avec les nouvelles technologies. L'analyse qualitative des données trouvées dans les transcriptions s'appuie sur l'analyse de contenu de Bardin \cite{bardin2008analise}, qui permet d'identifier des catégories conceptuelles et descriptives issues des discussions entre les participants. Les catégories descriptives, élaborées lors de l'analyse des données, peuvent détacher des aspects importants pour la recherche sur les affects vécus et englober des problématiques non abordées dans le cadre théorique qui constituent les catégories conceptuelles.

\section{Considérations finales}

Les premiers résultats soulignent l'importance de comprendre comment les jeunes interagissent sur des plateformes comme Instagram et confirment le rôle des affects pour faire face à cette vie complexe gouvernée par les algorithmes, qui représente le colonialisme des données \cite{couldry2019costs}. Les rapports recueillis lors des cinq groupes de discussion \cite{abreu2009grupos} confirment les limites du contrôle de l'utilisateur sur sa visualité, même pour ceux qui connaissent déjà l'esthétique \og instagrammable\fg{}%
\footnote{Catégorie identifiée pendant les groupes de discussion \cite{abreu2009grupos} à partir de l'analyse de contenu de Bardin \cite{bardin2008analise}. Les participants aux groupes de discussion \cite{abreu2009grupos} ont qualifié \og esthétique instagrammable\fg{} les caractéristiques esthétiques de la plateforme Instagram, telles que son cadrage, ses méthodes de navigation, son affichage et d'autres aspects.}
et possèdent des compétences numériques minimales, qui s'avèrent insuffisantes face à la complexité imposée par la médiation algorithmique de cette plateforme.

L'esthétique d'Instagram façonne la visualité de l'utilisateur, provoquant un phénomène de dissociation qui explique en grande partie les processus de fragmentation et la hausse des taux de détresse psychologique chez les jeunes. L'observation des affects dans les relations médiatisées par les algorithmes des réseaux sociaux nous aide à mieux comprendre ces processus, caractéristiques de cet isolement sous l'égide de l'hyperconnectivité. Ce phénomène est intimement lié à la désinformation et à l'exclusion, en raison du manque de contrôle dans leurs expériences sur Instagram. Un phénomène a été exprimé par des affects de méfiance, d'étrangeté et d'anxiété. Il faut noter que certains outils proposés par Instagram se sont avérés inefficaces pour protéger l'utilisateur et contribuent à le configurer esthétiquement, ce qui caractérise un nouveau modèle de pouvoir.

Les expériences proposées lors des groupes de discussion \cite{abreu2009grupos} ont démontré à quel point la dimension affective est déterminante dans une éducation au visuel numérique. Cela inclut également l'établissement de protocoles de protection dans ces espaces contrôlés et influencés par des algorithmes. L'analyse des données indique que les expériences photographiques programmées peuvent représenter une stratégie de protection pour ces jeunes, sachant que cette première phase de la recherche a révélé d'importantes limites à cet égard.

Nous sommes convaincus qu'une éducation pour le visuel numérique permet d'élargir la littératie numérique, en intégrant la dimension affective de l'individu, au-delà du contrôle colonisateur des algorithmes. Dans cette optique, il est essentiel d'intégrer une culture éthique et juridique à l'éducation au visuel numérique. Une éducation qui privilégie la dimension affective et l'expérience photographique, reconnaissant le pouvoir des images et contribuant à la sensibilisation et au développement de l'individu dès l'enfance. Il est donc pertinent de proposer cette approche dans différents contextes académiques, scolaires et universitaires.

\section*{Financement}

Ce projet est développé dans le cadre de la recherche postdoctorale en psychologie sociale de l'auteur, sous la supervision de la professeure Edna Lúcia Tinoco Ponciano, financé par la FAPERJ (Fondation pour la Recherche de l'État de Rio de Janeiro), et a été approuvé par le Comité d'Éthique de l'Université d'État de Rio de Janeiro – CAAE: 89037624.0.0000.5282. La recherche a débuté en mars 2025 et est développée en collaboration avec DERA-UERJ.

\printbibliography

\end{document}


1.Selon Joanna Zylinska [14], professeure de philosophie des médias et de pratiques numériques critiques au King's College de Londres, le terme « post-photographie » s'est imposé au cours de la dernière décennie pour décrire la nouvelle réalité du médium photographique. Ce terme a émergé du débat sur les similitudes et les différences entre la photographie analogique e numérique, ainsi que sur le rôle des « amateurs » dans la production photographique. Le terme « post-photographie » a gagné en visibilité en 2011, suite à un manifeste de Joan Fontcuberta, publié dans le journal La Vanguardia, qui annonçait une nouvelle ère pour la photographie, qui a laissé dans le passé certaines caractéristiques de ce médium dans la construction de la réalité et de la mémoire. Un terme qui acquiert une nouvelle pertinence avec les systèmes algorithmiques génératifs et les nouvelles formes de production d'images, comme ce qu'on appelle déjà la du « Photo Prompt » ou « Photographie du Prompt ». Il est important de noter que Zylinska utilise le terme « After Photography» [14], proche de la dimension affective, liée à l'idée fondamentale de notre projet.
Nous vous recommandons quelques lectures sur le sujet: FONTCUBERTA, Joan. The Post-Photographic Condition. Kerber, 2016; FONTCUBERTA, Joan. Welcome to the post-photographic era, 2017, disponible à: <https://www.barcelona.cat/metropolis/en/contents/welcome-the-post-photographic-era>; BEJENARU, Matei et al.. Post-Photography/Post-Fotografia. Vector, 2015;  MOREIRAS, Camila. Joan Fontcuberta: post-photography and thespectral image of saturation, Journal of Spanish Cultural Studies, 2017, disponible à: https://doi.org/10.1080/14636204.2016.1274496;  SCHWAB, Michael. Image Automation: Post-Conceptual Post-Photography and the Deconstruction of the Photographic Image. Thèse présentée en vue de l'obtention partielle du doctorat en philosophie au Royal College of Art, 2008; FERREIRA, Ângela. Reframing colonial narratives: Notes about Post-photography; disponible à: <https://ler.letras.up.pt/uploads/ficheiros/19312.pdf>; SHRIYAL, S. et al. The Evolution of Post-Photography Art in an Era Dominated by AI Image Synthesis. ShodhKosh: Journal of Visual and Performing Arts, 7(4s), 16–24. doi:   10.29121/shodhkosh.v7.i4s.2026.7476, 2026; GARCÍA SEDANO, M. Una revisión del concepto de postfotografía. Imágenes contra el poder desde la red. Liño, 21(21), 125–132. https://doi.org/10.17811/li.21.2015.125-132, 2015;  Understanding Post-Photography: exploiting the iconosphere; disponible à: <https://aestheticsofphotography.com/understanding-post-photography-exploiting-the-iconosphere/>.

2. Catégorie identifiée pendant les groupes de discussion [1] à partir de l´analyse de contenu de Bardin [2].
Les participants aux groupes de discussion [1] ont qualifié « esthétique instagrammable » les caractéristiques esthétiques de la plateforme Instagram, telles que son cadrage, ses méthodes de navigation, son affichage et d'autres aspects.

